We consider the following integral with parameters~$\l$,~$\f$:
\begin{align*}
I(\l,\f):=\int_{\R}{\rm e}^{{\rm i}\l S(\m,\f)}\c(\m,\l,\f){\rm d}\m
\end{align*}
and its decay rate with respect to $\l\gg 1$

If we freeze the parameter $\f$, then the stationary phase theorem just implies that if $\m\mapsto S(\m,\f)$ is a Morse function (that is, all critical points are non-degenerate in the sense that $\pa_{\m}^2S\neq 0$ there), the optimal decay rate of $I$ is \smash{$\l^{-\frac{1}{2}}$}. To obtain an improved decay, here we consider the following two scenarios:
\begin{itemize}\itemsep=0pt
\item If we assume that $\c$ vanishes with order $2\n$ at all the critical points of $S$ in addition, then the optimal decay rate is improved to be \smash{$\l^{-\frac{1}{2}-\n}$}.
\item If the symbol $\c$ itself decays like $(1+\l)^{-\n}$, then the decay order becomes \smash{$\l^{-\frac{1}{2}-\n}$}.
\end{itemize}
The following proposition interpolates the above two situations in a uniform way with respect to the parameter $\f$.
A typical example of such a phase function is $S(\m,\f)=(\m-\f)^2$. This simple case could simplify the reading of the proof below.

\begin{Proposition}\label{Stphasemovecrit}
Let $\n\geq 0$ and $c, C>0$. Suppose that $S\in C^{\infty}\bigl([0,1]^2;\R\bigr)$ and $\c(\cdot,\l,\cdot)\in C^{\infty}(\R\times [0,1])$ satisfy
\begin{align*}
&\supp \c(\cdot,\l,\f)\subset [0,1],\qquad |\pa_{\m}^{\a}\c(\m,\l,\f)|\leq C_{\a}\m^{2\n-\a}(1+\l \f\m)^{-\n},\\
&
\big|\pa_{\m}^2S(\m,\f)\big|\geq c,\qquad |\pa_{\m}\pa_{\f}S(\m,\f)|\geq C%\leq C
\end{align*}
uniformly in $\m\in [0,1]$, $\f\in [0,1]$ and $\l\gg 1$. Suppose that $S(\cdot,\f)$ has a unique critical point~${\m_0(\f)}$ which is smooth with respect to $\f$ and that $\m_0(0)=0$. Then
 \smash{$|I(\l,\f)|\lesssim \l^{-\n-\frac{1}{2}}$} uniformly in $\l\gg 1$ and $\f\in [0,1]$.
\end{Proposition}

\begin{Remark}\quad
\begin{itemize}\itemsep=0pt
\item[(1)] Without the assumption $\m_0(0)=0$, we obtain \smash{$I(\l,\f)\lesssim \l^{-\frac{1}{2}}$} only. The additional decay~$\l^{-\n}$ comes from the factor $(1+\l \f\m)^{-\n}$ in the assumption of $\c$ and the fact that $\m_0(0)=0$ as is mentioned before.

\item[(2)] We observe that the critical point $\m_0(\f)$ of $S$ is close to $0$ if $\f\sim 0$. There we take advantage of the assumption that the symbol $\c$ vanishes at zero with order $2\n$ like $|\c(\m,\l,\f)|\lesssim \m^{2\n}$. On the other hand, if $\f$ is away from zero, then the critical point $\m_0(\f)$ is also away from zero. In this case, $\c$ itself decays like $\l^{-\n}$ near the critical point (due to $|\f|\gtrsim 1$ and~${|\m_0(\f)|\gtrsim 1}$), which makes us to prove the improved decay. The difficulty here is to deal with the case where $\f$ is small enough but depending on $\l$.

\item[(3)]
The critical point of $S(\cdot,\f)$ is unique due to the condition $\big|\pa_{\m}^2S(\m,\f)\big|\geq c$ in this case.
\end{itemize}
\end{Remark}

\begin{proof}[Proof of Proposition~\ref{Stphasemovecrit}]
We may assume \smash{$\supp \c(\cdot,\l,\f)\subset \bigl[\l^{-\frac{1}{2}},1\bigr]$}. In fact, we set
\begin{align*}
\c'(\m,\l,\f)=\chi\bigl(\l^{\frac{1}{2}}\m\bigr)\c(\m,\l,\f),
\end{align*}
where $\chi\in C^{\infty}([0,\infty),\R)$ satisfies $\chi(\m)=1$ on $\m\leq 1$ and $\chi(\m)=0$ for $\m\geq 2$. Using the bound~${|\c(\m,\l,\f)|\lesssim \m^{2\n}}$, we have
\begin{align*}
\left|\int_{\R}{\rm e}^{{\rm i}\l S(\m,\f)}\c'(\m,\l,\f){\rm d}\m\right|\lesssim \l^{-\n}\left|\int_{\R}\chi\bigl(\l^{\frac{1}{2}}\m\bigr){\rm d}\m\right|\lesssim \l^{-\n-\frac{1}{2}}.
\end{align*}
Hence, we can replace $\c$ by $\c-\c'$, where we note that $\c-\c'$ satisfies
\begin{align*}
\supp (\c-\c')(\cdot,\l,\f)\subset\bigl[\l^{-\frac{1}{2}},1\bigr],\qquad |\pa_{\m}^{\a}(\c-\c')(\m,\l,\f)|\leq C_{\a}\m^{2\n-\a}(1+\l \f\m)^{-\n}.
\end{align*}
In the following, we assume \smash{$\supp \c(\cdot,\l,\f)\subset \bigl[\l^{-\frac{1}{2}},1\bigr]$}.

 $(i)$ First, we consider the case $\f\in [0,1]$ satisfying \smash{$\m_0(\f)\leq 2^{-1}\l^{-\frac{1}{2}}$}.
Since the signature of~$\pa_{\m}^2S(\m,\f)$ does not change for \smash{$\m\in \bigl[\l^{-\frac{1}{2}},1\bigr]$} and $\f\in [0,1]$, we have
\[
|\pa_{\m}S(\m,\f)|=\left|\int_{\m_0(\f)}^{\m}\pa_{\m}^2S(\m',\f){\rm d}\m'\right|=\int_{\m_0(\f)}^{\m}\big|\pa_{\m}^2S(\m',\f)\big|{\rm d}\m'\geq c(\m-\m_0(\f))\geq 2^{-1}c\m
\]
 for \smash{$\m\in \bigl[\l^{-\frac{1}{2}},1\bigr]$}. By integrating by parts and using \smash{$\supp \c(\cdot,\l,\f)\subset \bigl[\l^{-\frac{1}{2}},1\bigr]$}, we have
\begin{align*}
|I(\l,\f)|=\left|(-{\rm i}\l)^{-N}\int_{\R}{\rm e}^{{\rm i}\l S(\m,\f)} L^N\c(\m,\l,\f){\rm d}\m\right|\leq \l^{-N}\int_{\l^{-\frac{1}{2}}}^1|L^N\c(\m,\l,\f)|{\rm d}\m,
\end{align*}
where $L=\pa_{\m}\circ (\pa_{\m}S)(\m,\f)^{-1}$ and we take $N>\n+1$. We observe from the assumption and Lemma \ref{Leibnizrule} that $\big|L^N\c(\m,\l,\f)\big|\lesssim \m^{2\n-2N}$, which leads to
\begin{align*}
|I(\l,\f)|\lesssim \l^{-N}\int_{\l^{-\frac{1}{2}}}^1\m^{2\n-2N}{\rm d}\m\lesssim \l^{-\n-\frac{1}{2}}.
\end{align*}

 $(ii)$ Next, we consider the case $\f\in [0,1]$ satisfying \smash{$\m_0(\f)\in \bigl[ 2^{-1}\l^{-\frac{1}{2}}, 1\bigr]$}.
Differentiating $(\pa_{\m}S)(\m_0(\f),\f)=0$ with respect to $\f$, we have
\begin{align*}
|\m_0'(\f)|=\bigl|-(\pa_{\m}\pa_{\f}S)(\m_0(\f),\f)/\pa_{\m}^2S(\m_0(\f),\f)\bigr|\sim 1,
\end{align*}
and hence $\m_0(\f)\sim \f$ for $\f\in [0,1]$ due to $\m_0(0)=0$. Now the assumption \smash{$\m_0(\f)\geq 2^{-1}\l^{-\frac{1}{2}}$} yields \smash{$\f\gtrsim \l^{-\frac{1}{2}}$}.

By scaling, we have
\begin{align*}
I(\l,\f)=\int_{\R}{\rm e}^{{\rm i}\l S(\m,\f)}\c(\m,\l,\f){\rm d}\m=\m_0(\f)\l^{-\n}\int_{\R}{\rm e}^{{\rm i}\l \m_0(\f)^2S_{\f}(\m)}\c_{\f}(\m,\l){\rm d}\m,
\end{align*}
where we set
\begin{align*}
S_{\f}(\m)=\m_0(\f)^{-2}S(\m_0(\f)\m,\f),\qquad \c_{\f}(\m,\l)=\l^{\n}\c(\m_0(\f)\m,\l,\f).
\end{align*}
We note that $\m_0(\f)$ is a unique critical point of $S(\m,\f)$.
Let $\g_1,\g_2,\g_3\in C^{\infty}(\R;[0,1])$ such that~${\g_1+\g_2+\g_3=1}$, $\supp \g_1\subset \bigl(-\infty,\frac{3}{4}\bigr]$, $\supp \g_2\subset \bigl[\frac{1}{2},\frac{3}{2}\bigr]$ and $\supp \g_3\subset \bigl[\frac{5}{4},\infty\bigr)$. We write
\begin{align*}
I(\l,\f)
=\m_0(\f)\l^{-\n}\sum_{j=1}^3\int_{\R}{\rm e}^{{\rm i}\l \m_0(\f)^2S_{\f}(\m)}\c_{\f}(\m,\l)\g_j(\m){\rm d}\m
=\m_0(\f)\l^{-\n}\sum_{j=1}^3I_j(\l,\f).
\end{align*}

First, we deal with $I_2(\l,\f)$.
Now we see $|\pa_{\m}S(\m,\f)|\geq c|\m-\m_0(\f)|$ for $\m\in [0,1]$ and $\f\in [0,1]$. Then we have
\begin{align*}
|\pa_{\m}S_{\f}(\m)|=\m_0(\f)^{-1}\cdot |(\pa_{\m}S)(\m_0(\f)\m,\f)|\geq c|\m-1|
\end{align*}
for $\m\in \supp \c_{\f}(\cdot,\l)$ and $\f\in (0,1]$.
Moreover, for $\a\in\mathbb{N}$ and $\a'\in \mathbb{N}\setminus\{0\}$, we have
\begin{gather*}
|\pa_{\m}^{\a}\c_{\f}(\m,\l)|=\l^{\n} \m_0(\f)^{\a}|(\pa_{\m}^{\a}\c)(\m_0(\f)\m,\l,\f)|\\
\phantom{|\pa_{\m}^{\a}\c_{\f}(\m,\l)|}{}\lesssim \l^{\n}\m_0(\f)^{\a} (\m_0(\f) \m)^{2\n-\a}\bigl(1+\l \m_0(\f)^2 \m\bigr)^{-\n}\lesssim \m^{\n-\a},\\
\supp (\c_{\f}(\cdot,\l))\subset \big\{\m_0(\f)^{-1}\l^{-\frac{1}{2}}\leq \m\leq \m_0(\f)^{-1}\big\},\qquad \big|\pa_{\m}^{1+\a'}S_{\f}(\m)\big|\lesssim |\m_0(\f)|^{\a'-1}
\end{gather*}
Thus the stationary phase theorem (see Lemma \ref{stphaselem} with $k=2$, $j=0$ and $x_0=1$) implies
\begin{align*}
\m_0(\f)\l^{-\n}|I_2(\l,\f)|
\lesssim \m_0(\f)\l^{-\n}\cdot \bigl(\l \m_0(\f)^2\bigr)^{-\frac{1}{2}}=\l^{-\frac{1}{2}-\n},
\end{align*}
where we use $\l \m_0(\f)^2\gtrsim 1$.

Next, we consider the estimate of $I_{1}(\l,\f)$. We note that $|\pa_{\m}S_{\f}(\m)|\gtrsim 1$ for $\m\in \supp \g_1$.
Then the non-stationary phase theorem (see Lemma \ref{singularasym} with $\m=\n$) implies
\begin{align*}
\m_0(\f)\l^{-\n}|I_1(\l,\f)|\lesssim \m_0(\f)\l^{-\n}\cdot \bigl(\l \m_0(\f)^2\bigr)^{-\n-1}\lesssim \f\l^{-\n}\cdot \bigl(\l \f^2\bigr)^{-\frac{1}{2}}=\l^{-\frac{1}{2}-\n},
\end{align*}
where we use $\l \m_0(\f)^2\gtrsim 1$.

Finally, we consider the term $I_{3}(\l,\f)$. We note that $|\pa_{\m}S_{\f}(\m)|\gtrsim (|\m|+1)$ for $\m\in \supp \g_3$. Then, by integrating by parts $N(>1)$ times, we have
\begin{align*}
\m_0(\f)\l^{-\n}|I_3(\l,\f)|\lesssim \m_0(\f)\l^{-\n}\cdot \bigl(\l \m_0(\f)^2\bigr)^{-N}\lesssim \m_0(\f)\l^{-\n}\cdot \bigl(\l \m_0(\f)^2\bigr)^{-\frac{1}{2}}=\l^{-\frac{1}{2}-\n},
\end{align*}
where we use $\l \m_0(\f)^2\gtrsim 1$. This completes the proof.
\end{proof}
